\textbf{Время:} 01:45

\textbf{Тезис:} кривизна траектории даёт импульс, а ионизационные потери помогают определить вид частицы.

\textbf{На экране:} слева две траектории и знак магнитного поля; затем формула для p\_T; справа график ионизационных потерь и подпись о сравнении с расчётными значениями для π, K и p.

\textbf{Текст:} В TPC есть магнитное поле, поэтому заряженная частица движется по искривлённой траектории. Сравним две дуги слева. При одинаковых поле и модуле заряда чем меньше изгиб, тем больше поперечный импульс. Поперечный означает направленный перпендикулярно пучку; его обозначают p-T.

Формула под рисунком выражает эту связь через радиус кривизны r и магнитное поле B. Направление изгиба позволяет определить знак заряда. Если дополнительно учесть наклон пространственной траектории к пучку, получаем полный вектор импульса.

Но одинаковый импульс могут иметь разные частицы. Посмотрим на правый график: по горизонтали отложен импульс с учётом знака заряда, по вертикали — потери энергии на ионизацию на единицу пути. Пионы, каоны и протоны образуют разные полосы. Их обозначения — π, K и p.

Для каждого трека сравниваем измеренные потери с расчётными значениями для разных частиц. Полосы местами перекрываются, поэтому только по этому признаку частицы удаётся различить не во всём диапазоне импульсов.

\textbf{Переход:} Добавим независимое измерение — время пролёта частицы. Вместе с импульсом оно даст оценку массы.

\textbf{Справка:} В практической формуле p\_T измеряется в ГэВ/c, B — в теслах, r — в метрах; Q/e — заряд в единицах элементарного. Для π± модуль Q/e равен единице. На оригинальной оси графика указано p/q; это обозначение относится к заряду частицы и не совпадает с нашим относительным импульсом пары q. Кривизна даёт p\_T/|Q|, а наклон трека — продольную компоненту. Anderson et al., NIM A 499, 659–678 (2003), §5.7; А. Краева (STAR), Universe 10, 188 (2024), рис. 2. График относится к Au+Au при 3 ГэВ, 0–80\%.
